\section{FEM 与全局刚度矩阵基础}

\begin{frame}{从结构问题到矩阵方程}
\begin{block}{有限元离散后的典型形式}
$$K u = f$$
\end{block}
\begin{itemize}
  \item $K$：global stiffness matrix，描述整体结构刚度。
  \item $u$：全局自由度位移向量。
  \item $f$：外载荷、边界条件处理后的右端项。
  \item assembly 的目标：把每个单元的局部刚度矩阵 $K_e$ 累加到 $K$。
\end{itemize}
\source{standard FEM formulation; docs/requirements/cpu-parallel-stiffness-assembly-design.md}
\end{frame}
\note{
新手理解时不要一上来陷入代码。有限元最终是把连续体问题变成代数方程组。这里的项目只研究其中一个关键步骤：如何构造 K。K 构造完成以后，还要施加边界条件并求解 Ku=f，但那不是当前主线。
}

\begin{frame}{什么是 element stiffness matrix}
\begin{columns}[T,onlytextwidth]
\begin{column}{0.50\textwidth}
\textbf{局部视角}
\begin{itemize}
  \item 每个 element 有自己的节点。
  \item 每个节点有若干 DOF。
  \item 单元内部形成小矩阵 $K_e$。
\end{itemize}
\end{column}
\begin{column}{0.46\textwidth}
\textbf{全局视角}
\begin{itemize}
  \item 不同 element 共享全局节点。
  \item 局部 DOF 要映射到全局 DOF。
  \item 多个 $K_e$ 贡献到同一个 $K$。
\end{itemize}
\end{column}
\end{columns}
\takeaway{$K_e$ 是局部刚度，$K$ 是所有局部刚度按连接关系累加后的整体刚度。}
\end{frame}
\note{
可以把每个 element 想象成局部小结构。它有自己的节点编号和局部自由度编号。assembly 的关键动作就是 local-to-global mapping：找到局部自由度在全局向量里的位置，然后把 Ke(i,j) 加到 K(global_i, global_j)。
}

\begin{frame}{DOF：自由度是矩阵维度的来源}
\begin{itemize}
  \item DOF 是 degree of freedom，即节点上未知量的编号。
  \item 当前 C++ 主线默认每节点 \metric{3} 个位移自由度：$u_x,u_y,u_z$。
  \item 如果节点数是 $N$，全局位移向量长度约为 $3N$。
  \item WindTurbineHub 约 \metric{228384} nodes，因此 DOFs 约 \metric{685152}。
\end{itemize}
\source{README.md; results/2026-05-11-symbolic-numeric/symbolic\_numeric\_eval\_report.md}
\end{frame}
\note{
导师提到 section/closure 时，本质上也在谈“网格对象和数据自由度如何关联”。当前项目没做通用 Section，因为目前固定每个节点 3 个位移 DOF，直接映射已经足够支撑当前 benchmark。
}

\begin{frame}{Tet4、Hex8 与 Abaqus 单元名}
\begin{tabularx}{\textwidth}{lXX}
\toprule
术语 & 含义 & 当前项目角色 \\
\midrule
\code{Tet4} & 4 节点四面体单元 & 当前真实工程算例主单元 \\
\code{Hex8} & 8 节点六面体单元 & 支持规则网格和 C3D8 线弹性验证 \\
\code{C3D4} & Abaqus 里的 3D 4-node tetrahedral element & \code{3d-WindTurbineHub.inp} 中的真实输入类型 \\
\code{C3D8} & Abaqus 里的 3D 8-node brick element & 对应 Hex8/C3D8 物理刚度路径 \\
\bottomrule
\end{tabularx}
\takeaway{\code{Tet4}/\code{Hex8} 是代码单元类型，\code{C3D4}/\code{C3D8} 是 Abaqus 输入中的对应名字。}
\end{frame}
\note{
不同软件对单元的命名不同。项目代码里说 Tet4/Hex8，Abaqus .inp 文件里写 C3D4/C3D8。当前主汇报 case 仍是 WindHub 的 Tet4/C3D4，但代码里的正式刚度模型已经支持 Tet4 和 Hex8 两条物理路径。
}

\begin{frame}{mesh topology 与 stiffness model 是分层配置}
\begin{columns}[T,onlytextwidth]
\begin{column}{0.48\textwidth}
\textbf{mesh topology}
\begin{itemize}
  \item 单元有几个节点。
  \item 节点如何连接成 element。
  \item 决定 DOF 映射和稀疏 pattern。
\end{itemize}
\end{column}
\begin{column}{0.48\textwidth}
\textbf{stiffness model}
\begin{itemize}
  \item 局部矩阵 $K_e$ 怎么计算。
  \item \code{linear\_elastic\_solid}：当前正式模型。
  \item 按 topology 分派到 Tet4 或 Hex8 物理公式。
\end{itemize}
\end{column}
\end{columns}
\takeaway{两层在软件架构中分开处理，但真实物理计算要求 topology 与 stiffness model 合法匹配。}
\source{src/assembly/element\_kernels.cpp; docs/requirements/cpu-parallel-stiffness-assembly-design.md}
\end{frame}
\note{
这里要避免把“两层配置”讲成“任意独立”。Topology 决定节点数、DOF 映射、CSR pattern 和 scatter 位置；stiffness model 决定每个单元的 $K_e$ 怎么算。真实 FEM 里它们有合法组合约束：Tet4/C3D4 走 constant-strain Tet4，Hex8/C3D8 走 2x2x2 Gauss full integration。
}

\begin{frame}{\code{legacy\_synthetic}/\code{simplified} 的边界}
\begin{itemize}
  \item \code{simplified} 现在只是 \code{legacy\_synthetic} 的历史别名。
  \item 必须显式允许后才可用于小型 smoke/provenance。
  \item 它生成确定性合成 $K_e$，便于隔离 assembly strategy。
  \item 它不再支撑正式 benchmark 或 report-facing 结论。
\end{itemize}
\takeaway{当前正式结论只使用 \code{linear\_elastic\_solid}，不要用 \code{simplified} 代表真实 FEM。}
\source{README.md; docs/context/current-knowledge-boundary.md}
\end{frame}
\note{
如果听众问 simplified 是不是还能说明问题，回答要分清范围：它仍保留 element-to-global assembly、CSR scatter 和冲突结构，所以可以解释早期 assembly-overhead provenance；但它不代表真实材料、几何和积分成本，不能再作为当前正式 benchmark 结论。
}

\begin{frame}{\code{linear\_elastic\_solid} 如何匹配 topology}
\begin{itemize}
  \item Canonical CLI：\code{--stiffness-model linear\_elastic\_solid}。
  \item Tet4/C3D4：计算体积、形函数梯度、$B$、$D$，形成 $K_e = B^T D B \cdot V$。
  \item Hex8/C3D8：使用 2x2x2 Gauss full integration。
  \item \code{physics\_tet4} 只是 Tet4/C3D4-only 历史入口。
  \item $K_e$ 仍按 scatter 位置累加到全局 CSR values。
\end{itemize}
\takeaway{真实物理路径不是任意组合，而是由 \code{linear\_elastic\_solid} 根据 element type 分派。}
\source{src/assembly/element\_kernels.cpp; README.md}
\end{frame}
\note{
这里不需要推导完整弹性力学公式。理解重点是当前代码的正式入口已经不是“simplified vs physics_tet4”二选一，而是以 linear_elastic_solid 作为统一模型，再按 element type 调到 Tet4 或 Hex8 的物理局部刚度实现。physics_tet4 这个名字只作为历史 Tet4-only alias 保留。
}

\begin{frame}{为什么全局刚度矩阵是稀疏矩阵}
\begin{itemize}
  \item 每个单元只连接少数节点。
  \item 一个 DOF 只和相邻单元中的 DOF 有直接刚度耦合。
  \item 绝大多数远距离 DOF 之间没有直接矩阵条目。
  \item 因此 $K$ 的非零项远少于 dense matrix 的 $n^2$ 项。
\end{itemize}
\takeaway{稀疏性来自有限元网格的局部连接关系。}
\end{frame}
\note{
如果用 dense matrix 保存 WindTurbineHub 的 K，维度约 685k x 685k，内存会不可接受。CSR 这类稀疏格式只保存非零结构和非零值，是工程上必须使用的表示。
}
